\chapter{Passiv}
\label{cha:Passiv}

Das Passiv drückt aus, dass die Handlung im Vordergrund steht und der Handelnde unbekannt oder unwichtig ist.

\section{Vorgangspassiv}
\label{sec:Vorgangspassiv}

Bildung: werden + Partizip II

\begin{itemize}
    \item Das Haus wird gebaut. (gerade)
    \item Die Rechnung wurde bezahlt. (vergangen)
    \item Das Buch wird gelesen werden. (Zukunft)
\end{itemize}

\section{Zustandspassiv}
\label{sec:Zustandspassiv}

Bildung: sein + Partizip II

\begin{itemize}
    \item Das Fenster ist geöffnet. (= Zustand)
    \item Die Tür war geschlossen.
    \item Die Arbeit ist gemacht.
\end{itemize}

\section{Passiversatzformen}
\label{sec:Passiversatz}

\subsection{man-Konstruktion}
\begin{itemize}
    \item Das Fahrrad wird repariert. $\rightarrow$ Man repariert das Fahrrad.
    \item Hier wird Deutsch gesprochen. $\rightarrow$ Man spricht hier Deutsch.
\end{itemize}

\subsection{sein + zu + Infinitiv}
\begin{itemize}
    \item Das Haus wird gebaut. $\rightarrow$ Das Haus ist zu bauen. (= man muss)
    \item Die Aufgabe ist zu lösen. (= man muss)
\end{itemize}

\subsection{lassen + sich}
\begin{itemize}
    \item Das Problem wird gelöst. $\rightarrow$ Das Problem lässt sich lösen.
    \item Der Text ist zu verstehen. $\rightarrow$ Der Text lässt sich verstehen.
\end{itemize}

\section{Übungen zum Passiv}
\label{sec:PassivUebungen}

\begin{enumerate}
    \item Jemand repariert das Fahrrad. $\rightarrow$ Passiv \quad \textit{(Lösung: Das Fahrrad wird repariert.)}
    \item Sie baut ein Haus. $\rightarrow$ Passiv \quad \textit{(Lösung: Ein Haus wird gebaut.)}
    \item Man hat das Fenster geöffnet. $\rightarrow$ Zustandspassiv \quad \textit{(Lösung: Das Fenster ist geöffnet.)}
    \item Jemand hat die Tür geschlossen. $\rightarrow$ Zustandspassiv \quad \textit{(Lösung: Die Tür ist geschlossen.)}
    \item Das Fahrrad wird repariert. $\rightarrow$ mit „man" \quad \textit{(Lösung: Man repariert das Fahrrad.)}
    \item Das Haus wird gebaut. $\rightarrow$ mit „sein+zu" \quad \textit{(Lösung: Das Haus ist zu bauen.)}
    \item Ich bringe das Auto zur Werkstatt. $\rightarrow$ mit „lassen" \quad \textit{(Lösung: Ich lasse das Auto reparieren.)}
\end{enumerate}

\chapterend